\chapter{Membilang sampai 20}\label{ch:g1:counting}

Berapa banyak kelereng? Membilang menjawab pertanyaan pertama dalam
matematika. Pada bab ini kita membilang kumpulan benda yang kecil,
menulis bilangan $0$ sampai $20$, lalu membandingkan: siapa yang punya
lebih banyak?

\section{Membilang benda}

\begin{method}[Cara membilang]\label{met:g1:counting:how}
\begin{enumerate}
  \item Sentuh setiap benda satu kali sambil menyebut bilangan
        berurutan: satu, dua, tiga, \dots;
  \item jangan sampai ada benda yang terlewat, jangan menyentuh satu
        benda dua kali (geser benda yang sudah dibilang, atau beri
        tanda silang);
  \item bilangan \emph{terakhir} yang kamu sebut adalah banyaknya
        benda.
\end{enumerate}
Urutan menyentuh tidak berpengaruh: dibilang dari kiri ke kanan atau
dari kanan ke kiri, $7$ kelereng tetap $7$ kelereng.
\end{method}

\begin{omfigure}
\begin{tikzpicture}[scale=0.75]
  \foreach \i/\n in {0/1,1/2,2/3,3/4,4/5} {
    \fill[omDef] (\i*1.1,0) circle (0.33);
    \node[above, font=\small] at (\i*1.1,0.45) {\n};
  }
  \node[right, font=\small] at (5.2,0) {lima kelereng: ada $5$};
\end{tikzpicture}

{\small Membilang: satu bilangan untuk setiap kelereng, dan bilangan
terakhir adalah jawabannya.}
\end{omfigure}

\begin{definition}[Nol]\label{def:g1:counting:zero}
Ketika tidak ada yang bisa dibilang --- kotak kosong --- bilangannya
adalah \emph{nol}\index{nol}, ditulis $0$.
\end{definition}

\section{Bilangan sampai 20}

\begin{example}[Membaca dan menulis]\label{ex:g1:counting:writing}
Bilangan $0$ sampai $20$, berurutan:
\[
0,\ 1,\ 2,\ 3,\ 4,\ 5,\ 6,\ 7,\ 8,\ 9,\ 10,
\]
\[
11,\ 12,\ 13,\ 14,\ 15,\ 16,\ 17,\ 18,\ 19,\ 20 .
\]
Setelah $9$, bilangan ditulis dengan \emph{dua} angka: $13$ adalah satu
puluhan dan $3$ satuan; $20$ adalah dua puluhan. (Puluhan mendapat
gilirannya pada \cref{ch:g1:tens}.)
\end{example}

\begin{omfigure}
\begin{tikzpicture}[scale=0.62]
  \draw[-Stealth] (-0.4,0) -- (20.8,0);
  \foreach \i in {0,...,20} {
    \draw (\i,-0.12) -- (\i,0.12);
    \node[below, font=\footnotesize] at (\i,-0.15) {\i};
  }
  \fill[omThm] (7,0) circle (2.6pt);
  \fill[omMeth] (13,0) circle (2.6pt);
\end{tikzpicture}

{\small Garis bilangan dari $0$ sampai $20$: setiap bilangan punya
tempatnya. Berjalan ke kanan, bilangan makin besar. Titik menandai $7$
dan $13$.}
\end{omfigure}

\begin{example}[Tepat sebelum, tepat sesudah]\label{ex:g1:counting:neighbors}
Pada garis bilangan, setiap bilangan punya dua tetangga: tepat sebelum
$14$ ada $13$, tepat sesudahnya ada $15$. Membilang mundur --- $10$,
$9$, $8$, \dots\ --- berarti berjalan ke kiri.
\end{example}

\section{Membandingkan}

\begin{method}[Siapa yang lebih banyak?]\label{met:g1:counting:compare}
Untuk membandingkan dua kumpulan, buatlah pasangan: satu benda dari
masing-masing kumpulan, berdampingan.
\begin{enumerate}
  \item Jika ada benda pada salah satu kumpulan yang tidak kebagian
        pasangan, kumpulan itu \emph{lebih banyak};
  \item jika tidak ada yang tersisa di kedua sisi, banyaknya
        \emph{sama} pada keduanya.
\end{enumerate}
Dengan bilangan, kita tidak perlu memasangkan: yang letaknya lebih ke
kanan pada garis bilangan itulah yang lebih besar --- memasangkan dan
garis bilangan selalu sepakat. Kita menulis $12 > 9$ ($12$ lebih
banyak) atau $9 < 12$ ($9$ lebih sedikit).
\end{method}

\begin{omfigure}
\begin{tikzpicture}[scale=0.62]
  \foreach \i in {0,...,6} \fill[omDef] (\i*1.1,1.3) circle (0.3);
  \foreach \i in {0,...,4} \fill[omThm] (\i*1.1,0) circle (0.3);
  \foreach \i in {0,...,4}
    \draw[omExo, thick] (\i*1.1,0.3) -- (\i*1.1,1.0);
  \draw[omExo, thick] (5.5,1.3) circle (0.52);
  \draw[omExo, thick] (6.6,1.3) circle (0.52);
  \node[right, font=\small] at (7.5,0.65)
    {$7 > 5$: dua kelereng biru tidak punya pasangan};
\end{tikzpicture}

{\small Membandingkan dengan memasangkan: $7$ kelereng biru melawan
$5$ kelereng merah --- dua kelereng biru (yang dilingkari) tersisa,
jadi $7$ lebih banyak daripada $5$.}
\end{omfigure}

\begin{example}[Pertama, kedua, ketiga]\label{ex:g1:counting:ordinal}
Dalam lomba lari, para pelari sampai berurutan: yang \emph{pertama},
yang \emph{kedua}, yang \emph{ketiga}, yang \emph{keempat}\,\dots\
Kata-kata ini menunjukkan \emph{urutan}, bukan banyaknya: pelari
ketiga adalah satu orang pelari, di tempat nomor $3$.
\end{example}

\section{Latihan}

\begin{exercise}[$\star$]\label{exo:g1:counting:1}
Bilanglah huruf pada nama depanmu. Bilanglah jari pada dua tangan.
Bilanglah kaki dua ekor kucing.
\end{exercise}

\begin{exercise}[$\star$]\label{exo:g1:counting:2}
Tulislah dengan angka: tujuh; \ dua belas; \ dua puluh; \ \omterm{def:g1:counting:zero}{nol}.
Tulislah dengan kata: $4$; \ $15$; \ $18$.
\end{exercise}

\begin{exercise}[$\star$]\label{exo:g1:counting:3}
Lanjutkan bilangannya: $8,\ 9,\ 10,\ ?,\ ?,\ ?$. Bilanglah mundur:
$15,\ 14,\ 13,\ ?,\ ?,\ ?$.
\end{exercise}

\begin{exercise}[$\star$]\label{exo:g1:counting:4}
Bilangan mana yang tepat sebelum $10$? Tepat sesudah $16$? Di antara
$12$ dan $14$?
\end{exercise}

\begin{exercise}[$\star$]\label{exo:g1:counting:5}
Salin dan lengkapi dengan $<$ atau $>$:
\[
6 \;?\; 9, \qquad
14 \;?\; 8, \qquad
11 \;?\; 13, \qquad
20 \;?\; 19 .
\]
\end{exercise}

\begin{exercise}[$\star$]\label{exo:g1:counting:6}
Gambarlah $6$ lingkaran dan $8$ tanda silang. Pasangkan keduanya: mana
yang lebih banyak, dan berapa lebihnya?
\end{exercise}

\begin{exercise}[$\star$]\label{exo:g1:counting:7}
Urutkan dari yang terkecil ke yang terbesar: $12$; \ $7$; \ $17$; \ $2$; \ $10$.
\end{exercise}

\begin{exercise}[$\star$]\label{exo:g1:counting:8}
Lima anak berbaris: Ani, Budi, Citra, Dani, Eka. Siapa yang kedua?
Siapa yang keempat? Di urutan ke berapa Dani?
\end{exercise}

\begin{exercise}[$\star\star$]\label{exo:g1:counting:9}
Aku sebuah bilangan antara $10$ dan $20$. Menurut garis bilangan, aku
tiga langkah sesudah $13$. Siapakah aku? Dan siapa yang tiga langkah
\emph{sebelum} $13$?

\end{exercise}
